\section{Introduction}

The assignment for this course was to build and deploy a distributed Guestbook application,
starting from local containerization and working towards a fully automated cloud-native
deployment on OpenShift. The idea was to go through the whole pipeline myself, from writing
a Containerfile to pushing images and getting everything running in a real cluster.

The application itself has four parts working together. There is an Nginx-based frontend
that the user talks to, a Go backend that handles the actual logic, a PostgreSQL database
for storing guestbook entries, and a Redis cache that sits in front of the database to
speed up reads. When a user loads the page, the frontend forwards the request to the
backend, which checks Redis first before hitting PostgreSQL. This setup is a good
reflection of how real production services are built.

The original plan was to deploy directly to the shared OpenShift cluster at
\texttt{devops24.cloud2.se}, configure ArgoCD there, and wire everything together with
GitHub Actions. Part way through the course that changed and we moved to the OpenShift
Developer Sandbox instead. That shift ended up teaching me more than the original plan
would have, because the Sandbox has strict platform restrictions that forced me to really
understand what each piece of the system actually needs to function.

\subsection{Understanding the Repository}

Before writing any Containerfiles I spent time reading through the repository to understand
what the backend was doing. The backend is written in Go and uses Go modules, so the first
things I checked were \texttt{go.mod} and \texttt{go.sum} to see what dependencies were
involved and whether they could be downloaded cleanly. I then ran the backend locally,
hitting the \texttt{/health}, \texttt{/api/entries}, and \texttt{/api/stats} endpoints
manually to confirm it behaved as expected. I also checked that the environment variables
for both PostgreSQL and Redis were being read correctly, since those connections are central
to everything else working.
